\chapter{Thiết kế và triển khai hệ thống}
\label{chap:implementation}

\section{Mục tiêu triển khai}

Hệ thống được xây dựng để đáp ứng bốn yêu cầu: sinh mẫu có thể tiếp tục sau gián đoạn,
đo chi phí nhất quán, phát lại nhiều policy trên cùng dữ liệu và lưu đủ provenance để
kiểm tra claim. Mã nguồn Python được tổ chức thành package \code{adaptive\_compute},
tách logic sinh GPU khỏi phân tích CPU. Cách tách này cho phép sửa phân tích hoặc thêm
bảng thống kê mà không phải chạy lại mô hình, miễn là quy tắc parse và generation
identity không thay đổi.

\begin{figure}[htbp]
  \centering
  \fbox{\begin{minipage}{0.91\textwidth}
    \centering
    \textbf{Frozen dataset manifest} $\rightarrow$
    \textbf{Prompt + deterministic seed} $\rightarrow$
    \textbf{Qwen3 generation}\\[0.35cm]
    $\downarrow$ append-only checkpoint và generation cache\\[0.35cm]
    \textbf{Strict parser} $\rightarrow$
    \textbf{Valid-only aggregation} $\rightarrow$
    \textbf{Fixed/adaptive policy replay}\\[0.35cm]
    $\downarrow$\\[-0.1cm]
    \textbf{Cost accounting + statistical reports + figures}
  \end{minipage}}
  \caption{Luồng dữ liệu chính của hệ thống thực nghiệm.}
  \label{fig:system-pipeline}
\end{figure}

\section{Tổ chức mã nguồn}

Các module được phân trách nhiệm như sau:

\begin{itemize}
  \item \code{config.py}, dataset adapters và split modules định nghĩa cấu hình, tải
  dữ liệu và xác minh manifest;
  \item \code{prompts.py}, module adapter Transformers và
  \code{experiment\_runner.py} xây prompt, sinh mẫu và ghi provenance;
  \item \code{answer\_parser.py}, \code{evaluate.py} và \code{aggregation.py} parse,
  chuẩn hóa, chấm đúng/sai và tổng hợp cấu hình;
  \item \code{generation\_cache.py} cùng checkpoint JSONL hỗ trợ tái sử dụng và
  tiếp tục sau gián đoạn;
  \item \code{controller/} xây dữ liệu tuần tự, kiểm tra training gate, đóng băng và
  phát lại policy;
  \item \code{analysis/} tạo oracle, khoảng tin cậy, kiểm định ghép cặp, bảng Pareto
  và biểu đồ;
  \item \code{tests/} kiểm tra unit và integration bằng fake adapter, không yêu cầu
  CUDA.
\end{itemize}

Command-line interface là lớp mỏng gọi các module trên. Thực nghiệm GPU và phân tích
CPU có lệnh riêng, nhờ đó người dùng có thể kiểm tra artifact trước khi chuyển giai
đoạn.

\section{Adapter mô hình và phép đo}

\code{TransformersModelAdapter} tải tokenizer và causal language model theo model ID
cùng revision đã chốt. Mặc định thiết bị là CUDA và kiểu dữ liệu là bfloat16. Adapter
từ chối tự động rơi về CPU nếu người dùng không bật rõ \code{allow\_cpu}; cơ chế này
tránh vô tình chạy một job nhiều giờ trên CPU.

Prompt được đưa qua chat template của Qwen3 với thinking mode bật. Khi temperature
lớn hơn 0, adapter bật sampling với top-$p$ đã cấu hình; seed được đặt cho PyTorch và
CUDA trước mỗi lần sinh. Thời gian được đo sau khi đồng bộ CUDA ở hai đầu. Số prompt
token lấy từ tensor input và số token sinh lấy từ phần output nằm sau prompt, không
ước lượng từ chuỗi đã decode. Finish reason là \code{length} khi số token chạm cap,
\code{stop} khi gặp EOS, hoặc trạng thái chuẩn hóa khác khi không xác định được.

Cấu hình chính dùng Qwen3-0.6B, temperature 0,7, top-$p$ 0,95, thinking mode, CUDA và
bfloat16. Revision mô hình đã resolve được lưu trong manifest để tránh model ID trỏ
tới nội dung khác ở lần chạy sau.

\section{Định danh thí nghiệm và manifest bất biến}

Experiment ID được tạo từ model, dataset, số bài, global seed và hash của toàn bộ
generation identity. Identity gồm model revision, dataset revision, danh sách problem
ID, các giá trị $k,L$, số repetition, decoding settings, prompt template, parser và
seed policy. Manifest còn lưu hash split, package versions, device, dtype và số lần
sinh dự kiến.

Khi chạy mới, manifest được tạo theo chế độ độc quyền. Khi resume, chương trình so
sánh từng trường ảnh hưởng kết quả; chỉ một khác biệt về prompt hash, parser, seed,
revision hoặc action space cũng làm resume thất bại. Manifest vì vậy không chỉ là log
mô tả mà là rào chắn chống trộn hai thí nghiệm không tương thích.

\section{Checkpoint và generation cache}

Mỗi generation được append ngay vào \code{raw\_samples.jsonl}. Khóa checkpoint là

\begin{equation}
  (\text{experiment ID},\text{problem ID},L,\text{sample ID},
  \text{repetition},\text{model revision},\text{prompt hash}).
\end{equation}

$k$ cố ý vắng mặt vì breadth được dẫn xuất từ tiền tố của shared pool. Khi resume,
hệ thống lập chỉ mục bản ghi mới nhất cho mỗi khóa. Một dòng cuối bị ghi dở do tiến
trình dừng đột ngột có thể được bỏ qua an toàn; lỗi JSON ở giữa file được xem là hỏng
checkpoint và làm chương trình dừng. Bản ghi lỗi generation được giữ là outcome hoàn
chỉnh trừ khi người dùng yêu cầu retry rõ ràng.

Generation cache có identity chứa mọi setting ảnh hưởng nội dung sinh nhưng cũng loại
$k$. Cache chỉ được dùng khi identity khớp; sau đó kết quả vẫn đi qua parser và cost
accounting của experiment hiện tại. Cơ chế này đặc biệt quan trọng trên RTX 4060 vì
một lần chạy có thể kéo dài và bị gián đoạn, nhưng không được đánh đổi bằng việc tái
sử dụng mẫu sai cấu hình.

\section{Parser, tổng hợp và artifact trung gian}

Sau sinh, mỗi bản ghi lưu raw response, đáp án reasoning nếu có, final response answer,
failure type, truncation stage, token, latency, seed và audit flag. Parser name được
version hóa để thay đổi quy tắc parse không âm thầm làm thay đổi kết quả cũ.

\code{aggregation.py} nhận đúng tiền tố $k$ mẫu, bỏ phiếu trên đáp án cuối hợp lệ và
tạo một dòng cho mỗi bài/repetition/profile. Dòng này chứa dự đoán majority, agreement,
valid fraction, correctness, số lần truncate/parse failure, ngân sách $kL$, token thật
sự, tổng prompt token và latency. Bảng \code{config\_results.csv} là cầu nối giữa raw
generation và mọi phân tích cấp policy.

Các artifact chính và vai trò được tóm tắt ở Bảng~\ref{tab:artifacts}.

\begin{center}
\begin{minipage}{0.96\textwidth}
  \captionof{table}{Artifact phục vụ tái lập và kiểm toán.}
  \label{tab:artifacts}
  \centering
  \begin{tabular}{p{0.34\textwidth}p{0.56\textwidth}}
    \toprule
    Artifact & Nội dung \\
    \midrule
    \code{experiment\_manifest.json} & Cấu hình bất biến, hash, revision, seed,
    problem ID và số generation dự kiến. \\
    \code{raw\_samples.jsonl} & Một dòng trên mỗi mẫu với response, parse, lỗi,
    token, latency và identity. \\
    \code{config\_results.csv} & Một dòng trên mỗi bài và $(k,L)$ với accuracy,
    agreement và chi phí. \\
    \code{policy\_traces.csv} & Quyết định tuần tự, lý do dừng và các mẫu policy
    thực sự đã mua. \\
    \code{paired\_comparisons.csv} & McNemar, Holm và chênh lệch metric ghép cặp. \\
    Báo cáo/đồ thị & Wilson interval, Pareto, diagnostics và trường hợp minh họa. \\
    \bottomrule
  \end{tabular}
\end{minipage}
\end{center}

\section{Phát lại policy và kiểm tra tính tương đương}

Policy replay đọc raw pool, yêu cầu repetition ID bằng 0 cho protocol hiện tại, từ
chối generation error, duplicate key hoặc pool thiếu. Với mỗi bài và $L$, sample ID
phải đúng dãy $0,1,2,3$ và question/ground truth phải nhất quán. Prompt-token count
cũng phải giống nhau giữa các pool của một bài.

Ngưỡng tertile được tính riêng từ Train. Tại mỗi quyết định, replay chỉ cung cấp cho
policy tiền tố đã trả phí. Trace lưu policy ID, policy type, partition, $L$, breadth đã
chọn, decision reason, prompt token và thresholds. Một kiểm tra endpoint xác nhận rằng
replay profile cố định tái tạo chính xác kết quả tổng hợp trực tiếp; nếu không, phân
tích bị dừng thay vì tiếp tục với bảng không nhất quán.

\section{Kiểm thử và môi trường tái lập}

Test suite hiện có 203 test và đã chạy thành công sau đợt refactor package. Các test
bao phủ parser, aggregation, seed và checkpoint identity, nested sample pool, dataset
manifest, statistical routines, controller data, policy freeze và policy replay. Fake
model adapter được dùng cho integration test để kiểm tra toàn pipeline mà không tải
mô hình hoặc yêu cầu GPU.

Môi trường thực nghiệm đã được kiểm tra với Python 3.11.15, PyTorch 2.9.1+cu128, CUDA
12.8 và NVIDIA GeForce RTX 4060. Dependency được khóa trong requirements và hướng dẫn
tái lập ghi rõ lệnh kiểm tra môi trường, chạy test, sinh GPU, resume và phân tích CPU.
Kết quả này chứng minh phần mềm hoạt động trong môi trường hiện tại; nó không thay thế
việc báo cáo revision, manifest và raw artifact của từng experiment.

\section{Giới hạn triển khai hiện tại}

Hệ thống đã triển khai non-thinking comparator bằng
\code{enable\_thinking=false}, với chế độ sinh được đưa vào manifest và cache identity.
Hệ thống cũng đã triển khai \code{force\_answer\_v2}: pha đầu nhận $L-R$ token;
nếu chạm giới hạn, pipeline đóng thinking block, thêm cue trả lời và dành tối đa
$R$ token còn lại mà không vượt tổng cap $L$. Strategy, version, reserve, trạng thái
forced transition và số control token chèn được lưu trong manifest/raw/cache identity.
Backend cố định không cung cấp native numeric thinking-budget. Learned joint
controller không vượt training gate nên không được dùng cho claim chính. Latency chỉ
có ý nghĩa trong cùng phần cứng và phiên chạy;
actual tokens được ưu tiên làm cost có khả năng so sánh hơn. Cuối cùng, deterministic
seed giúp tái lập pool nhưng một seed không đo được độ nhạy theo ngẫu nhiên. Các mục
Các giới hạn còn lại được ghi nhận thay vì trình bày như chức năng đã hoàn thành.
